\documentclass{amsart}
% For dealing with references we use the comment environment
\usepackage{verbatim}
\newenvironment{reference}{\comment}{\endcomment}
%\newenvironment{reference}{}{}
\newenvironment{slogan}{\comment}{\endcomment}
\newenvironment{history}{\comment}{\endcomment}

% For commutative diagrams we use Xy-pic
\usepackage[all]{xy}

% We use 2cell for 2-commutative diagrams.
\xyoption{2cell}
\UseAllTwocells

% We use multicol for the list of chapters between chapters
\usepackage{multicol}

% This is generally recommended for better output
\usepackage{lmodern}
\usepackage[T1]{fontenc}

% For cross-file-references

% Package for hypertext links:
\usepackage{hyperref}

% For any local file, say "hello.tex" you want to link to please

% Theorem environments.
%
\theoremstyle{plain}
\newtheorem{theorem}[subsection]{Theorem}
\newtheorem{proposition}[subsection]{Proposition}
\newtheorem{lemma}[subsection]{Lemma}

\theoremstyle{definition}
\newtheorem{definition}[subsection]{Definition}
\newtheorem{example}[subsection]{Example}
\newtheorem{exercise}[subsection]{Exercise}
\newtheorem{situation}[subsection]{Situation}

\theoremstyle{remark}
\newtheorem{remark}[subsection]{Remark}
\newtheorem{remarks}[subsection]{Remarks}

\numberwithin{equation}{subsection}

% Macros
%
\def\lim{\mathop{\mathrm{lim}}\nolimits}
\def\colim{\mathop{\mathrm{colim}}\nolimits}
\def\Spec{\mathop{\mathrm{Spec}}}
\def\Hom{\mathop{\mathrm{Hom}}\nolimits}
\def\Ext{\mathop{\mathrm{Ext}}\nolimits}
\def\SheafHom{\mathop{\mathcal{H}\!\mathit{om}}\nolimits}
\def\SheafExt{\mathop{\mathcal{E}\!\mathit{xt}}\nolimits}
\def\Sch{\mathit{Sch}}
\def\Mor{\mathop{\mathrm{Mor}}\nolimits}
\def\Ob{\mathop{\mathrm{Ob}}\nolimits}
\def\Sh{\mathop{\mathit{Sh}}\nolimits}
\def\NL{\mathop{N\!L}\nolimits}
\def\CH{\mathop{\mathrm{CH}}\nolimits}
\def\proetale{{pro\text{-}\acute{e}tale}}
\def\etale{{\acute{e}tale}}
\def\QCoh{\mathit{QCoh}}
\def\Ker{\mathop{\mathrm{Ker}}}
\def\Im{\mathop{\mathrm{Im}}}
\def\Coker{\mathop{\mathrm{Coker}}}
\def\Coim{\mathop{\mathrm{Coim}}}

% Boxtimes
%
\DeclareMathSymbol{\boxtimes}{\mathbin}{AMSa}{"02}

%
% Macros for moduli stacks/spaces
%
\def\QCohstack{\mathcal{QC}\!\mathit{oh}}
\def\Cohstack{\mathcal{C}\!\mathit{oh}}
\def\Spacesstack{\mathcal{S}\!\mathit{paces}}
\def\Quotfunctor{\mathrm{Quot}}
\def\Hilbfunctor{\mathrm{Hilb}}
\def\Curvesstack{\mathcal{C}\!\mathit{urves}}
\def\Polarizedstack{\mathcal{P}\!\mathit{olarized}}
\def\Complexesstack{\mathcal{C}\!\mathit{omplexes}}
% \Pic is the operator that assigns to X its picard group, usage \Pic(X)
% \Picardstack_{X/B} denotes the Picard stack of X over B
% \Picardfunctor_{X/B} denotes the Picard functor of X over B
\def\Pic{\mathop{\mathrm{Pic}}\nolimits}
\def\Picardstack{\mathcal{P}\!\mathit{ic}}
\def\Picardfunctor{\mathrm{Pic}}
\def\Deformationcategory{\mathcal{D}\!\mathit{ef}}

\setlength{\emergencystretch}{3em}
\begin{document}
\title[Stacks Project, Tag 0698]{490 --- Pseudo-coherent morphisms are syntomic local on the source}
\maketitle
\noindent Copyright (C) 2005--2025 Johan de Jong. Stacks Project authors. Source: \href{https://stacks.math.columbia.edu/tag/0698}{Tag 0698}.

\noindent Editorial difficulty note (not author text). Difficulty H2: The proof must descend pseudo-coherence through a syntomic cover rather than only through a fixed flat base change. It lifts the cover across a polynomial embedding, chooses a principal neighbourhood where the lift is faithfully flat, and translates pseudo-coherence to the defining ideal for faithfully flat module descent..

\noindent Editorial provenance note (not author text). History: excerpt from revision \texttt{a0444\allowbreak 6e57e\allowbreak c1fbc\allowbreak 252a8\allowbreak 71afc\allowbreak ec775\allowbreak 2fb28\allowbreak 07b14}, more-morphisms.tex, lines 17544--17618. Reference presentation was adapted; original-author context and core support blocks were retained or added. The mathematical wording is unchanged. Licensed under GNU FDL 1.2 or later; no invariant sections or cover texts. See ../licenses/COPYING. Context blocks reproduce author definitions and supporting results; they are not additional counted problems.

\begin{definition}
\label{definition-pseudo-coherent}
A morphism of schemes $f : X \to S$ is called {\it pseudo-coherent}
if the equivalent conditions of
Lemma \ref{lemma-pseudo-coherent}
are satisfied. In this case we also say that $X$ is pseudo-coherent
over $S$.
\end{definition}

\begin{lemma}
\label{lemma-pseudo-coherent}
Let $f : X \to S$ be a morphism of schemes. The following are equivalent
\begin{enumerate}
\item there exist an affine open covering $S = \bigcup V_j$ and for each $j$
an affine open covering $f^{-1}(V_j) = \bigcup U_{ji}$ such that
$\mathcal{O}_S(V_j) \to \mathcal{O}_X(U_{ij})$ is a pseudo-coherent
ring map,
\item for every pair of affine opens $U \subset X$, $V \subset S$
such that $f(U) \subset V$ the ring map
$\mathcal{O}_S(V) \to \mathcal{O}_X(U)$ is pseudo-coherent, and
\item $f$ is locally of finite type and $\mathcal{O}_X$
is pseudo-coherent relative to $S$.
\end{enumerate}
\end{lemma}

\begin{lemma}
\label{lemma-pseudo-coherent-syntomic-local-source}
The property $\mathcal{P}(f) =$``$f$ is pseudo-coherent''
is syntomic local on the source.
\end{lemma}

\begin{proof}
We will use the criterion of
Descent, Lemma \href{https://stacks.math.columbia.edu/tag/036H}{lemma properties morphisms local source (Tag 036H)}
to prove this. It follows from
Lemmas \href{https://stacks.math.columbia.edu/tag/0695}{lemma flat finite presentation pseudo coherent (Tag 0695)} and
\href{https://stacks.math.columbia.edu/tag/0681}{lemma composition pseudo coherent (Tag 0681)}
that being pseudo-coherent is preserved under precomposing with
flat morphisms locally of finite presentation, in particular under
precomposing with syntomic morphisms (see
Morphisms, Lemmas \href{https://stacks.math.columbia.edu/tag/01UL}{lemma syntomic flat (Tag 01UL)} and
\href{https://stacks.math.columbia.edu/tag/01UK}{lemma syntomic locally finite presentation (Tag 01UK)}).
It is clear from
Definition \ref{definition-pseudo-coherent}
that being pseudo-coherent is
Zariski local on the source and target.
Hence, according to the aforementioned
Descent, Lemma \href{https://stacks.math.columbia.edu/tag/036H}{lemma properties morphisms local source (Tag 036H)}
it suffices to prove the following: Suppose $X' \to X \to Y$ are
morphisms of affine schemes with $X' \to X$ syntomic and $X' \to Y$
pseudo-coherent. Then $X \to Y$ is pseudo-coherent.
To see this, note that in any case $X \to Y$ is of finite presentation by
Descent, Lemma \href{https://stacks.math.columbia.edu/tag/02KK}{lemma flat finitely presented permanence algebra (Tag 02KK)}.
Choose a closed immersion $X \to \mathbf{A}^n_Y$. By
Algebra, Lemma \href{https://stacks.math.columbia.edu/tag/00T0}{lemma lift syntomic (Tag 00T0)}
we can find an affine open covering $X' = \bigcup_{i = 1, \ldots, n} X'_i$
and syntomic morphisms $W_i \to \mathbf{A}^n_Y$ lifting the morphisms
$X'_i \to X$, i.e., such that there are fibre product diagrams
$$
\xymatrix{
X'_i \ar[d] \ar[r] & W_i \ar[d] \\
X \ar[r] & \mathbf{A}^n_Y
}
$$
After replacing $X'$ by $\coprod X'_i$ and setting $W = \coprod W_i$
we obtain a fibre product diagram
$$
\xymatrix{
X' \ar[d] \ar[r] & W \ar[d]^h \\
X \ar[r] & \mathbf{A}^n_Y
}
$$
with $W \to \mathbf{A}^n_Y$ flat and of finite presentation and
$X' \to Y$ still pseudo-coherent. Since
$W \to \mathbf{A}^n_Y$ is open (see
Morphisms, Lemma \href{https://stacks.math.columbia.edu/tag/01UA}{lemma fppf open (Tag 01UA)})
and $X' \to X$ is surjective we can find
$f \in \Gamma(\mathbf{A}^n_Y, \mathcal{O})$ such that
$X \subset D(f) \subset \Im(h)$. Write
$Y = \Spec(R)$, $X = \Spec(A)$, $X' = \Spec(A')$
and $W = \Spec(B)$, $A = R[x_1, \ldots, x_n]/I$ and
$A' = B/IB$. Then $R \to A'$ is pseudo-coherent. Picture
$$
\xymatrix{
A' = B/IB & B \ar[l] \\
A = R[x_1, \ldots, x_n]/I \ar[u] & R[x_1, \ldots, x_n] \ar[l] \ar[u]
}
$$
By
Lemma \href{https://stacks.math.columbia.edu/tag/0697}{lemma quotient of flat finitely presented (Tag 0697)}
we see that $IB$ is pseudo-coherent as a $B$-module.
The ring map $R[x_1, \ldots, x_n]_f \to B_f$ is faithfully flat by
our choice of $f$ above. This implies that
$I_f \subset R[x_1, \ldots, x_n]_f$
is pseudo-coherent, see
More on Algebra, Lemma \href{https://stacks.math.columbia.edu/tag/068R}{lemma flat descent pseudo coherent (Tag 068R)}.
Applying
Lemma \href{https://stacks.math.columbia.edu/tag/0697}{lemma quotient of flat finitely presented (Tag 0697)}
one more time we see that $R \to A$ is pseudo-coherent.
\end{proof}
\end{document}
